\section{The damping-ratio coordinate and its qualifying rule}

\subsection{The qualifying rule}

A cell qualifies for the \(\chi\) axis only if its reported damping is the amplitude-damping \(\gamma\) of a damped-harmonic-oscillator (pole) response, the same quantity as \(\chi = \gamma/2\omega_0\). Transition linewidths, total-coherence (\(T_2\)) and pure-dephasing (\(T_\phi\), equivalently \(T_2^*\)) rates, and pressure-broadening coefficients are different physical quantities and are excluded however cleanly measured. This is why quantities that are not amplitude damping cannot share the coordinate: a dephasing rate and an oscillator damping are not silently placed on the same coordinate. The single worked example of Part~III is traced to its primary source, with the damping convention confirmed against the paper.

\subsection{The one-axis observation}

The regimes are one axis, not separate phenomena: under (\(\chi<1\)), critical (\(\chi=1\)), over (\(\chi>1\)). The critical point is a boundary systems cross, not a state a real system rests at, since \(\chi=1\) is where oscillation gives way to monotone decay, a threshold crossed rather than occupied.

A concrete over-damped anchor: the M-tilt phonon mode of CsPbBr\(_3\) at 350~K, from an atomistic-simulation study of the cubic-to-tetragonal transition \citep{fransson}, whose own equation for mode-coordinate autocorrelation gives two relaxation times at that temperature, \(\tau_S=0.31\)~ps and \(\tau_L=5.22\)~ps. Inverting the source's own relation (\(\tau_{S,L} = \tau/(1\mp\sqrt{1-(\omega_0\tau)^2})\), \(\tau=2/\Gamma\)) gives \(\Gamma = 1/\tau_S+1/\tau_L = 3.417\)~ps\(^{-1}\) and \(\omega_0 = 1/\sqrt{\tau_S\tau_L} = 0.786\)~ps\(^{-1}\), so \(\chi=\Gamma/2\omega_0=2.174\): genuinely over-damped, and the classifier returns monotone, stable, \(\chi>1\) for it with no oscillatory component. The round-trip through the source's own equation reproduces both published times to \(5\times10^{-15}\), and the regime survives a hypothetical 10\% perturbation on either time (Supplementary, \texttt{verify\_cspbbr3\_cell.py}). This is a simulated, not measured, linewidth, and is reported as such: the pass table above admits only directly measured population lifetimes, so this entry stays here rather than in it.

\subsection{The hinge}

The qualifying rule of Section~4.1 excludes dephasing, pressure broadening, and inhomogeneous width from the $\chi$ axis on the ground that they report different physical quantities. That exclusion is correct, and it raises a sharper question than it answers: given a measured linewidth, under what conditions may it legitimately be read as a mechanical $\chi$ at all? Part~III answers that question. The answer is not a second and third coordinate on which those excluded quantities are plotted (that would repeat the error the qualifying rule exists to prevent, placing physically distinct objects on a common-looking axis); it is a hierarchy of gates through which a linewidth must pass before Section~4.1's coordinate applies to it. The engine classifies whether a degeneracy is defective; the hierarchy classifies whether a linewidth may become a $\chi$. The two are the same epistemic move applied to two different objects, and neither reads its answer from the surface number.
